\section{Varieties with splitting cotangent sheaves}
\label{sec:4}

One relevant case that keeps appearing in the discussion of varieties with
projectively flat sheaf of reflexive differentials is the one where $\Omega^{[1]}_X$
decomposes as a direct sum of Weil divisorial sheaves. We discuss settings
where there are local or global decompositions of this form. The results of
this section will be used later in the proof of our main result, but might be of
independent interest. In relation to \cite{MR3030068}, the local results are
new, prompted by the appearance of singularities in our context, and the results
on varieties with globally split cotangent sheaf allow us to replace some of the
more involved arguments of \loccit\ by a clearer structure result.

\subsection{Varieties where $\Omega^{[1]}_X$ is locally split}
\label{sect:split}

Given a variety $X$ where $\Omega^{[1]}_X$ is projectively flat, for instance because
the flatness criterion of \cite[Th.\,1.6]{GKP20a} applies, the local description
of projectively flat sheaves, \cite[Prop.\,3.11]{GKP20a}, can often be used to
reduce to the case where $\Omega^{[1]}_X$ decomposes locally. We will show that this
already forces $X$ to have quotient singularities.

\begin{prop}\label{prop:icq}
Let $X$ be a normal complex space with klt singularities. Assume that the
sheaf $\Omega^{[1]}_X$ of reflexive differentials is of the form
\begin{equation}\label{eq:icq}
\Omega^{[1]}_X \cong \sL^{\oplus \dim X},
\end{equation}
where $\sL$ is reflexive of rank one. Then, $X$ has only isolated, cyclic
quotient singularities.
\end{prop}

We prove Proposition~\ref{prop:icq} in Section~\vref{ssec:popicq}.

\begin{rem}
If $X$ is a quasi-projective variety with klt singularities where $\Omega^{[1]}_X$
is of the form $\sL^{\oplus \dim X}$ for a Weil-divisorial sheaf $\sL$, then the proof
of Proposition~\ref{prop:icq} applies verbatim to show that $X$ is
Zariski-locally a quotient of a smooth variety. We do not use this fact
however.
\end{rem}

\subsection{Varieties where $\Omega^{[1]}_X$ is globally split}
\label{sec:4-2}

There are settings where $\Omega^{[1]}_X$ decomposes globally, for instance because
it is projectively flat and $X$ is known to be simply connected. This obviously
puts strong conditions on the global geometry of $X$; see \cite{MR1758581} for
further results in the smooth setup.

The following proposition, including the proof, is due to the referee. He also
suggested Theorem~\ref{thm:druel}, which is stronger than our original result
that included a pseudo-effectivity assumption. Our previous arguments were of
analytic nature and used a theorem of Demailly~\cite{Dem02}.

\begin{prop}\label{prop:druel}
Let $U$ be a connected complex manifold of dimension $n \geq 2$. Assume that
$\Omega^1_U \cong \sL^{\oplus n}$, where $\sL$ is invertible. Then,
$c_1(\sL) = 0 \in H^1 \bigl(U,\, \Omega^1_U \bigr)$.
\end{prop}
\begin{proof}
Consider the induced decomposition on cohomology
\[
H^1 \bigl(U,\, \Omega^1_U\bigr) \cong H^1\bigl(U,\, \sL \bigr)^{\oplus n}.
\]
For reasons that
will become apparent in a second, we are particularly interested in the
coordinate hyperplanes
\[
Z_i := \left\{ (\alpha_1, \dots, \alpha_n) \in H^1\bigl(U,\, \sL \bigr)^{\oplus n} \bigm| \alpha_i = 0
\right\} \subsetneq H^1 \bigl(U,\, \Omega^1_U \bigr).
\]
To see how the $Z_i$ come into play, restate the assumption as a decomposition
of the tangent bundle: $\sT_U \cong \sA_1 \oplus \cdots \oplus \sA_n$. The summands $\sA_\sbullet \cong \sL^*$ are
regular rank-one foliations on $U$, with normal bundles
\[
\sN_i = \sA_1 \oplus \cdots \oplus \what{\sA_i} \oplus \cdots \oplus \sA_n.
\]
In this setting, a theorem of Baum-Bott, \cite[Prop.\,3.3 \&
Cor.~3.4]{MR0261635} but see also \cite[Lem.\,3.1]{MR1760872}, describes the
first Chern class $c_1(\sN_i)$ as lying in the subspace
\[
Z_i = H^1 \bigl(U,\, \sN_i^* \bigr) \subsetneq H^1 \bigl(U,\, \Omega^1_U \bigr).
\]
But since $c_1(\sL)$ is proportional to any of the $c_1(\sN_\sbullet)$, we find that
$c_1(\sL)$ is contained in $Z_1 \cap\cdots\cap Z_n = \{0\}$. The assertion is thus shown.
\end{proof}

Applying Proposition~\ref{prop:druel} to the smooth locus of a klt variety, we
obtain the following characterisation of quasi-abelian varieties.

\begin{thm}\label{thm:druel}
Let $X$ be a normal, projective klt variety of dimension $n \geq 2$. Assume that
the sheaf $\Omega^{[1]}_X$ of reflexive differentials is of the form
$\Omega^{[1]}_X \cong \sL^{\oplus n}$, where $\sL$ is reflexive of rank one. Then, $X$ is
quasi-abelian.
\end{thm}
\begin{proof}
We have seen in Proposition~\ref{prop:druel} that the $\bQ$-Cartier sheaves $\sL$
are numerically trivial, and then so is $\omega_X \cong \sL^{[\otimes n]}$. Since $X$ is klt,
it follows from \cite[Prop.\,V.4.9]{Nakayama04} or \cite[Th.\,4.2]{MR2134273}
that $\omega_X$ and $\sL$ are torsion. As a consequence, we find a finite
quasi-étale cover $\gamma: \what{X} \to X$ such that the reflexive pull-back
$\gamma^{[*]} \sL$ is trivial. The space $\what{X}$ is klt and its tangent sheaf
$\sT_{\what{X}} \cong \left(\gamma^{[*]} \sL \right)^{\oplus n}$ is likewise trivial. By the
solution of the Lipman-Zariski conjecture for spaces with klt singularities,
\cite[Th.\,6]{GKKP11}, it follows that the covering space $\what{X}$ is
smooth. In other words, $\what{X}$ is a parallelisable projective manifold.
By a classic result of Wang, $\what{X}$ is an Abelian variety; see for example
\cite[\S 3.9, Lem.\,1]{MR1334091}.
\end{proof}

\subsection{Proof of Proposition~\ref*{prop:icq}}
\label{ssec:popicq}

The proof of Proposition~\ref{prop:icq} uses the following lemma that might be
of independent interest.

\begin{lem}\label{lem:glcs}
Let $X$ be a normal, complex space with only cyclic quotient singularities.
Let $S \subseteq X_{\sing}$ be any irreducible component of the singular locus. Then,
there exists a non-empty, open set $S^\circ \subseteq S$ such that every $x \in S^\circ$ admits an
open neighbourhood $U=U(x) \subseteq X$ of the form $U = W \times V$, where $W$ is smooth
and $V$ has an isolated quotient singularity.
\end{lem}
\begin{proof}
To begin, consider the special case where $V$ is a complex vector space,
$\rho : \bZ_m \to \GL(V)$ is a linear representation of a cyclic group and where
$X \subseteq V / \bZ_m$ is an open neighbourhood of $[\vec 0]$. Recall that $\rho$
decomposes into direct a sum of one-dimensional representations. Identifying
$\bZ_m$ with roots of unity, we can thus find linear coordinates on $V$ such
that multiplication with $\xi \in \bZ_m$ is given as
\[
\xi\cdot (v_1, \dots, v_m) = \bigl(\xi^{n_1}\cdot v_1, \dots, \xi^{n_m}\cdot v_m\bigr).
\]
An elementary computation in these coordinates shows that $S^\circ$ exists, and can
in fact be chosen to be dense in $S$.

Returning to the general case where $X$ is arbitrary, let $s \in S$ be any
point. By~assumption, there exists an open neighbourhood $U = U(x)$ and a
cyclic cover $\gamma : \wtilde U \to U$, say with group $G$. Shrinking $U$ and
passing to a subgroup, if need be, we may assume without loss of generality
that $\gamma$ is totally branched over $s$, that is $\gamma^{-1}(s) = \{ \wtilde s \}$.
The cyclic group thus acts on $\wtilde U$ and fixes $\wtilde s$. The group
$G$ clearly acts on the vector space $T_{\wtilde U}|_{\wtilde s}$ and
linearisation at the fixed point, \cite[Proof of Th.\,4]{MR0084174} or
\cite[Cor.\,2 on p.\,13]{MR782881}, shows the existence of a $G$-invariant open
neighbourhood $\wtilde U'$ of $\vec 0 \in T_{\wtilde U}|_{\wtilde s}$ and an
equivariant, biholomorphic map between $\wtilde U$ and $\wtilde U$. In other
words, we are in the situation discussed in the first paragraph of this proof,
where the claim has already been shown.
\end{proof}

\begin{proof}[Proof of Proposition~\ref{prop:icq}]
As a first step in the proof of Proposition~\ref{prop:icq}, we claim that $\sL$
is $\bQ$-Cartier. In fact, taking determinants of both sides in \eqref{eq:icq},
we obtain that $\sL^{[\otimes \dim X]} \cong \omega_X$, which is $\bQ$-Cartier by assumption.
Choose a minimal number $N \in \bN^+$ such that $\sL^{[\otimes N]}$ is locally free. The
reflexive tensor product $\bigl(\Omega^{[1]}_U \bigr)^{[\otimes N]}$ is then likewise
locally free.

As a second step, we show that $X$ has cyclic quotient singularities. In
fact, given any $x \in X$, we find an open neighbourhood $U=U(x)$ over which
$\sL^{[\otimes N]}$ is trivial. Chose a trivialisation and let $\gamma: \wtilde{U} \to U$ be
the associated index-one cover, which is cyclic of order $N$. The complex
space $\wtilde{U}$ has again klt singularities and
$\gamma^{[*]}\, \sL \cong\nobreak \sO_{\wtilde{U}}$. In particular, it follows that
$\Omega^{[1]}_{\wtilde{U}} \cong \gamma^{[*]} \Omega^{[1]}_U$ and $\sT_{\wtilde U}$ are both free.
The solution of the Lipman-Zariski conjecture for spaces with klt
singularities, \cite[Th.\,3.8 \& comments after Th.\,1.1]{Druel13a} or
\cite[Th.\,1.13]{KS18}, then asserts that $\wtilde{U}$ is smooth. We conclude
that $X$ has cyclic quotient singularities only.

As a third and last step, we show that the singularities of $X$ are isolated.
Assume to the contrary and let $S \subseteq X_{\sing}$ be a positive-dimensional,
irreducible component. We have seen in Lemma~\ref{lem:glcs} that there a
non-empty, open subset $S^\circ \subseteq S$ such that every $x \in S^\circ$ admits an open
neighbourhood $U = U(x) \subseteq X$ of product form $U = W \times V$, where $W$ is smooth
and $V$ has an isolated quotient singularity. In particular,
$\dim W = \dim S$ and
\[
\Omega^{[1]}_U \cong \Omega^1_W \oplus \Omega^{[1]}_V.
\]
In particular, the reflexive tensor power $\bigl(\Omega^{[1]}_U \bigr)^{[\otimes N]}$ is
written as a direct sum of reflexive sheaves with
$\bigl(\Omega^{[1]}_W \bigr)^{[\otimes (N-1)]} \otimes \Omega^{[1]}_V$ as one of its direct
summands. It follows that $\bigl(\Omega^{[1]}_U \bigr)^{[\otimes N]}$ is not locally
free and accordingly, neither is $\bigl(\Omega^{[1]}_X \bigr)^{[\otimes N]}$. This
contradicts the results obtained in the first paragraph of this proof, and
therefore finishes the proof of Proposition~\ref{prop:icq}.
\end{proof}
